% Generated from validated synthesis. Do not hand-edit.
\section*{Monthly Signals}
\addcontentsline{toc}{section}{Monthly Signals}
\smallskip
\noindent\textbf{Frontier Models \& Voice} model更新はdefault experience、API lifecycle、voiceという利用surfaceの更新と一体化した。 \hfill{\footnotesize p.~\pageref{pkg:frontier-models-may}}\par
\smallskip
\noindent\textbf{Agents \& Runtime} agentを長く動かすため、sandbox、approval、network policy、CLI／MCP統合が実行環境の主要設計項目になった。 \hfill{\footnotesize p.~\pageref{pkg:agents-runtime-may}}\par
\smallskip
\noindent\textbf{Inference \& Serving} KV cache、MoE、kernel、runtimeの各層が同時に更新され、推論最適化がstack全体の問題として見えるようになった。 \hfill{\footnotesize p.~\pageref{pkg:inference-serving-may}}\par
\smallskip
\noindent\textbf{Agent Safety \& Security} 攻撃面はpromptだけでなくworkspace、long-term memory、retrieval、monitoring policyへ広がった。 \hfill{\footnotesize p.~\pageref{pkg:agent-safety-security-may}}\par
\smallskip
\noindent\textbf{Capability \& Evaluation} benchmark外の研究成果とthird-party evaluation設計が、frontier capabilityをどう確かめるかという問題を前景化した。 \hfill{\footnotesize p.~\pageref{pkg:capability-evaluation-may} / p.~\pageref{pkg:paper-watch-may}}\par
\medskip
\begin{claimboundary}[Retrospective scope]
本号は2026年5月を後日確認可能になった一次情報も用いて再構成するRetrospective Specialである。Coverage Periodと制作時点を同一視せず、vendor / project / author claimの境界は各記事とTechnical Notesで明示する。
\end{claimboundary}
\medskip
\tableofcontents

\medskip
\subsection*{May chronology — 選択済みEvidenceで見る5月}
\addcontentsline{toc}{subsection}{May chronology}
\begingroup
\small
\setlength{\tabcolsep}{4pt}
\renewcommand{\arraystretch}{1.08}
\begin{tabularx}{\linewidth}{@{}p{0.12\linewidth}X@{}}
\toprule
日付 & 本号で追う動き \\
\midrule
5月5日 & GPT-5.5 Instant更新。default experienceの更新としてfrontier modelが現れる。\autocite{src-01a8b52eee81} \\
5月7日 & Realtime voice向け新モデル。modelと低遅延配信を一体で読む必要が強まる。\autocite{src-24a3c172de48} \\
5月8日 & Running Codex safely。sandbox・approval・network policyがagent運用のcontrol planeとして明示される。\autocite{src-66be12207714} \\
5月13日 & Codex on Windowsのsandbox設計。host OS固有のisolationが実装課題になる。\autocite{src-4a3442a5ca8c} \\
5月15日 & vLLM v0.21.0。May release seriesの中でserving stackが継続更新される。\autocite{src-80b8d67933c0} \\
5月20日 & OpenAIが離散幾何の予想を反証したmodel事例を公表。benchmark外のcapability評価が論点になる。\autocite{src-1c93d6c1c725} \\
5月28日 & AnthropicがClaude Opus 4.8を公表。coding・agentic／long-running work向けの更新として位置付ける。\autocite{src-f6aa679babd7} \\
5月29日 & FlashInfer v0.6.12とthird-party evaluation playbook。runtime最適化と評価方法が同じ月末に前景化する。\autocite{src-74a7052acd7a,src-8c84b8422093} \\
\bottomrule
\end{tabularx}
\endgroup
\smallskip
\noindent{\footnotesize\color{SurveyMuted}日付は本号の選択済み一次資料で日単位に確認できたEventのみを掲載する。月精度のindex項目は日付を補完していない。}

